\documentclass[letterpaper,10pt]{article}
\usepackage[T1]{fontenc}
\usepackage[utf8]{inputenc}
\usepackage[spanish]{babel}
\usepackage[margin=0.85in]{geometry}
\usepackage[hidelinks]{hyperref}
\usepackage{titlesec}

\setlength{\parindent}{0pt}
\pagestyle{empty}
\raggedbottom
\titleformat{\section}{\large\bfseries}{}{0pt}{}
\titleformat{\subsection}{\normalsize\bfseries}{}{0pt}{}
\titlespacing*{\section}{0pt}{1em}{0.4em}
\titlespacing*{\subsection}{0pt}{0.7em}{0.3em}

\newcommand{\cventry}[4]{%
  \noindent\parbox[t]{0.74\textwidth}{\textbf{#1}}%
  \hfill
  \parbox[t]{0.24\textwidth}{\raggedleft #2}\\
  \textit{#3}\\
  #4\par\vspace{0.45em}
}

\newcommand{\certentry}[4]{%
  \noindent\parbox[t]{0.74\textwidth}{\textbf{#1}}%
  \hfill
  \parbox[t]{0.24\textwidth}{\raggedleft #2}\\
  \textit{#3}\\
  \textbf{Verificación:} #4\par\vspace{0.45em}
}

\newcommand{\certentrydesc}[5]{%
  \noindent\parbox[t]{0.74\textwidth}{\textbf{#1}}%
  \hfill
  \parbox[t]{0.24\textwidth}{\raggedleft #2}\\
  \textit{#3}\\
  #4\par
  \textbf{Verificación:} #5\par\vspace{0.45em}
}

\begin{document}

% Información personal
\begin{center}
    {\LARGE \textbf{Patricio Acevedo Flores}} \\
    \vspace{0.3em}
    Estudiante de Ingeniería Civil Industrial con Diploma en Tecnologías de la Información \\
    Pontificia Universidad Católica de Chile \\
    \href{mailto:psacevedo@uc.cl}{psacevedo@uc.cl} \quad | \quad +569 3269 3473 \\
    \href{https://linkedin.com/in/patricio-andrés-acevedo-flores-63300b30b/}{linkedin.com/in/patricio-acevedo}
\end{center}

% Perfil profesional
\section*{Perfil Profesional}
Ingeniero Civil Industrial en formación en la Pontificia Universidad Católica de Chile, con diploma en Tecnologías de la Información. Me especializo en transformar problemas complejos en soluciones implementables mediante analítica, automatización, inteligencia artificial y ciberseguridad. He desarrollado proyectos aplicados para el sector financiero y de salud, incluyendo plataformas de ciberinteligencia, gestión de incidentes, análisis de riesgos y automatización de procesos con agentes IA. Complemento esta experiencia con docencia universitaria en programación, optimización y bases de datos, fortaleciendo comunicación técnica, liderazgo y trabajo colaborativo. Busco contribuir en entornos de alto impacto donde la ingeniería, los datos y la tecnología generen mejoras concretas y medibles.


% Sección de Experiencia
\section*{Experiencia}

\subsection*{Experiencia Laboral}

\cventry
  {Practicante, Experiencia del Paciente}
  {feb. 2026--actualidad}
  {Novo Nordisk Chile}
  {Análisis de recorridos de pacientes, análisis de datos y automatización de procesos con agentes de IA.}

\cventry
  {Proyecto Riesgos de la Industria Financiera}
  {nov. 2025--abr. 2026}
  {CiberLab UC}
  {Desarrollo de una plataforma de análisis y visualización de incidentes y riesgos de la industria financiera.}

\cventry
  {Proyecto Hermes CTI}
  {sep. 2025--abr. 2026}
  {Asociación de Bancos e Instituciones Financieras --- CiberLab UC}
  {Creación de una plataforma de ciberinteligencia y gestión de incidentes para la Asociación de Bancos.}

\cventry
  {Proyecto Simulador de Ciberdefensa}
  {2024--2025}
  {Ejército de Chile --- CiberLab UC}
  {Desarrollo de un simulador para la protección de infraestructuras críticas en maquetas a escala.}

\cventry
  {Proyecto Análisis y Clasificación de Malware}
  {2024--2025}
  {Ejército de Chile --- CiberLab UC}
  {Desarrollo y optimización de técnicas de análisis de malware en entornos simulados mediante inteligencia artificial.}

\subsection*{Experiencia Académica}

\cventry
  {Ayudante, Programación como Herramienta para la Ingeniería}
  {2025-2--2026-1}
  {Pontificia Universidad Católica de Chile}
  {Corrección de laboratorios y apoyo a estudiantes.}

\cventry
  {Ayudante, Optimización}
  {2025-2--2026-1}
  {Pontificia Universidad Católica de Chile}
  {Apoyo en clases, resolución de ejercicios y asistencia en evaluaciones.}

\cventry
  {Ayudante, Bases de Datos}
  {2025-2--2026-1}
  {Pontificia Universidad Católica de Chile}
  {Atención de consultas y laboratorios, revisión de entregas y apoyo en evaluaciones.}

\cventry
  {Ayudante, Microeconomía}
  {2025-1}
  {Pontificia Universidad Católica de Chile}
  {Apoyo docente en bienestar y coordinación de casos.}

\cventry
  {Ayudante, Programación Avanzada}
  {2023--2024}
  {Pontificia Universidad Católica de Chile}
  {Apoyo en estructuras de datos, algoritmos avanzados y aplicaciones en Python.}

\cventry
  {Ayudante, Narración Gráfica de No Ficción}
  {2023--2024}
  {Pontificia Universidad Católica de Chile}
  {Apoyo en proyectos visuales y narrativas gráficas para trabajos de no ficción.}

\cventry
  {Ayudante, Introducción a la Programación}
  {2023-2}
  {Pontificia Universidad Católica de Chile}
  {Clases de apoyo y tutorías sobre algoritmos y conceptos básicos de programación.}

\cventry
  {Ayudante, Curso Desafíos de la Ingeniería}
  {2022-1}
  {Pontificia Universidad Católica de Chile}
  {Talleres prácticos y orientación en la resolución de problemas interdisciplinarios.}


% Sección de Intercambio Académico
\section*{Intercambio Académico}
\cventry
  {Facultad de Derecho}
  {2024}
  {Universidad de Barcelona}
  {Cursos: Cibercriminología, Cooperación Penal Internacional, Crimen Organizado y Terrorismo, Estructura Social (análisis de políticas públicas en Europa) y Seguridad y Salud Laboral.}

% Sección de Certificaciones
\section*{Certificaciones}

\certentry
  {ML Practitioner Certificate}
  {feb. 2026}
  {Dataiku}
  {dndw5wmdw99e}

\certentry
  {Dataiku Core Designer}
  {feb. 2026}
  {Dataiku}
  {86e6k379opa8}

\certentry
  {Building with the Claude API}
  {feb. 2026}
  {Anthropic}
  {p6pcrm8fjkts}

\certentrydesc
  {Toma mejores decisiones basadas en datos: Power BI}
  {2025}
  {Santander X}
  {Curso enfocado en análisis de datos y visualización para la toma de decisiones.}
  {SX-2025-0727000032097}

\certentry
  {Cloud Security Basics}
  {2024}
  {University of Minnesota --- Coursera}
  {\href{https://coursera.org/verify/AW4925RKP2C}{AW4925RKP2C}}

\certentry
  {International Cyber Conflicts}
  {2024}
  {State University of New York --- Coursera}
  {\href{https://coursera.org/verify/YMF74CFSLLPR}{YMF74CFSLLPR}}

% Sección de Habilidades
\section*{Habilidades}
\textbf{Técnicas:} Python, Paquete Office, Análisis de datos, Programación avanzada, SQL, R y Kali Linux.\\
\textbf{Blandas:} Trabajo en equipo, adaptabilidad y liderazgo.\\
\textbf{Herramientas:} Power BI, Jupyter, herramientas de reuniones digitales y gestión de proyectos.

\end{document}